% SSU Data Dash 2026: team report template. Compile from this folder (Overleaf: upload the whole
% repository as a zip). `python run.py report` writes the results table and figure it includes.
% Keep the report to 4 pages, not counting the outside-data table and the appendix.
\documentclass[11pt]{article}
\usepackage[margin=2.2cm]{geometry}
\usepackage{graphicx,booktabs,hyperref}
\graphicspath{{../figures/}}
\setlength{\parskip}{4pt}

% Required: your team name and the full name of every team member.
\title{SSU Data Dash 2026: \textit{Team name}}
\author{\textit{Full name of every team member: First Last, First Last, \ldots}}
\date{}

\begin{document}
\maketitle

\noindent\begin{tabular}{@{}ll}
\textbf{Video (YouTube, unlisted is fine):} & \url{https://youtu.be/YOUR-VIDEO} \\
\textbf{Code and submission files (GitHub):} & \url{https://github.com/YOUR-TEAM/YOUR-REPO} \\
\end{tabular}

\section*{Summary}
% Three or four sentences: your approach, your validation result, and your main finding.
\textit{What you did, how well it worked when tested on later sales, and the one thing you learned.}

\section{Data and features}
% Which columns you used, what you built from them, and anything you dropped and why.
\textit{Your features for each track.}

\section{Models}
% One paragraph per track. Off-market has no list price; list-aware does. Explain how your two
% models differ and why.
\textit{Your off-market model and your list-aware model.}

\section{Validation}
% How you tested the models without using information from the future: which months you
% trained on, which you predicted, and how you made sure nothing leaked.
\textit{How you validated.}

\section{Results}
\begin{table}[h]
\centering
\input{../reports/results_table.tex}
\caption{Walk-forward validation on the listings. MdAPE: median absolute percentage error.
MAPE: mean absolute percentage error. Brackets: bootstrap 95\% interval.}
\end{table}
% Describe the numbers in words. What improved over the baselines, and by how much?
\textit{What the table shows.}

\section{Where the model fails}
\begin{figure}[h]
\centering
\includegraphics[width=0.85\linewidth]{error_by_price_range.png}
\caption{Absolute percentage error by sold-price decile.}
\end{figure}
% Add your own breakdowns: by month, property type, neighbourhood, size.
\textit{Where errors are largest, and why you think so.}

\section{Limitations}
\textit{What your model cannot do, and what you would try next.}

\section*{Outside data}
% Every outside source, or write "None". Data may only cover periods up to and including May 2026.
\begin{tabular}{lll}
\toprule
Source & URL & Periods covered \\
\midrule
\textit{None} & & \\
\bottomrule
\end{tabular}

\section*{Reproducing our results}
% The exact commands that produce your two submission files from a fresh copy of your repository.
\texttt{pip install -r requirements.txt}\\
\texttt{python run.py validate \&\& python run.py report \&\& python run.py predict \&\& python run.py check}

\end{document}
